\subsection{Założenia projektowe}
W odpowiedzi na zidentyfikowane ograniczenia klasycznego podejścia dwuetapowego, zaproponowano opracowanie systemu umożliwiającego bezpośrednie przekształcenie sygnału audio w wektorową reprezentację semantyczną, kompatybilną z przestrzenią popularnych modeli embeddingowych tekstu. Kluczowa innowacja polega na pominięciu etapu jawnej transkrypcji tekstowej, co pozwala wyeliminować związane z nim błędy oraz zredukować łączną złożoność obliczeniową w porównaniu z klasycznym podejściem (ASR, a następnie embedding tekstu). Nadrzędnym założeniem jest, że model zdolny do uchwycenia znaczenia wypowiedzi bezpośrednio z jej reprezentacji akustycznej poprawi skuteczność wyszukiwania w kontekście danych specjalistycznych.

Proponowane rozwiązanie wykorzystuje sekwencyjne reprezentacje akustyczne pozyskane z zaawansowanych modeli automatycznego rozpoznawania mowy (ASR). Sekwencja ta jest następnie agregowana do postaci pojedynczego wektora cech. Do agregacji zostanie zastosowany mechanizm uwagi (attention pooling), który adaptacyjnie wyodrębnia i przypisuje wyższe wagi najbardziej informatywnym fragmentom sygnału. W alternatywnym wariancie architektury, w celu bardziej zaawansowanego modelowania zależności sekwencyjnych, wykorzystany zostanie dekoder oparty na architekturze transformatora. Tak uzyskana reprezentacja akustyczna jest następnie odwzorowywana do docelowej przestrzeni wektorów cech semantycznych tekstu za pomocą warstwy projekcji liniowej.

Model zostanie wytrenowany na parach sygnału audio i odpowiadającego mu tekstu przy założeniu maksymalizwoania podobieństwa cosinusowego. Celem uczenia jest minimalizacja odległości semantycznej między embeddingiem audio a embeddingiem tekstowym referencyjnej wypowiedzi, tak aby wektory reprezentujące to samo znaczenie znajdowały się blisko siebie w przestrzeni docelowej. Umożliwi to bezpośrednie porównywanie embeddingów pochodzących z obu modalności – audio i tekstu.

Opracowane warianty architektury zostaną poddane analizie porównawczej z wykorzystaniem metody Centered Kernel Alignment (CKA)\cite{pmlr-v97-kornblith19a}, która umożliwia ocenę stopnia podobieństwa między wewnętrznymi reprezentacjami modeli uczenia maszynowego. Analiza ta pozwoli określić, czy różne mechanizmy agregacji sekwencji uczą się podobnych reprezentacji semantycznych oraz czy uzyskany model charakteryzuje się stabilnymi właściwościami generalizacyjnymi.

W ramach praktycznej walidacji systemu zostanie zbudowana wektorowa baza wiedzy zawierająca wygenerowane, badania medyczne, opisy leków oraz symulowane dane pacjentów. Baza ta posłuży do porównania skuteczności semantycznego wyszukiwania informacji przy użyciu zapytań głosowych przetworzonych na embeddingi bezpośrednio z sygnału audio oraz – jako punkt odniesienia – embeddingów powstałych na podstawie referencyjnej transkrypcji tekstowej. Następnie zostanie opracowany system typu Retrieval-Augmented Generation (RAG), który będzie wykorzystywał zarówno tekst wygenerowany przez moduł ASR, jak i dokumenty odnalezione za pomocą bezpośrednich embeddingów audio. Zintegrowany model językowy będzie odpowiedzialny za przeformułowanie zapytania użytkownika do bardziej precyzyjnej postaci, pobranie najbardziej adekwatnych dokumentów na podstawie embeddingów oraz wygenerowanie finalnej, umotywowanej źródłowo odpowiedzi. Pozwoli to ocenić, czy połączenie obu ścieżek przetwarzania – bezpośredniej (audio-embedding) i pośredniej (audio-transkrypcja) – zwiększa dokładność i spójność odpowiedzi w krytycznym kontekście danych medycznych.

Kompleksowa ocena systemu obejmie zarówno standardowe miary, takie jak WER (Word Error Rate) dla modułu ASR, jak i jakościową analizę różnic między przestrzeniami reprezentacji uzyskanych z różnych wariantów modeli. W celu specjalistycznej oceny w domenie medycznej zostanie opracowana ważona macierz błędów, która przypisze różne wagi (kary) potencjalnym zniekształceniom informacji, priorytetyzując błędy krytyczne dla bezpieczeństwa, takie jak błędne nazwy leków, niepoprawne dawki lub mylące nazwy chorób, zabiegów. Zastosowanie tej metryki pozwoli zweryfikować, czy projektowany system spełnia podstawowe wymagania w zakresie wierności informacji i poprawia jakość wyszukiwania w bazach specjalistycznych, przez co może być rozważany jako użyteczne narzędzie wspomagające w zastosowaniach klinicznych.

\subsection{Implementacja}
\subsection{Eksperymenty}
\subsection{Analiza wyników}
